% A finite matrix product state in mixed canonical form, with a two-site gate
% about to be applied: one step of TEBD. The triangles say which tensors are
% isometric -- canl to the left of the centre, canr to the right -- and the
% arrows on the bonds say which way the gauge runs. The centre is the one
% tensor that is not isometric, and it is also the one drawn as a diamond:
% shape and colour both single it out, which is the point, because the norm and
% every local expectation value can be read off it alone.
%
% Finite, so the chain simply stops. The outer bonds of the first and last
% tensors are one-dimensional and carry nothing, and a picture of a finite
% state should not put an object there.
%
% The physical legs leave each triangle at the bottom corner of its flat side,
% which is \tnlegs{l} to the left of the centre and \tnlegs{r} to the right:
% the two gauges then mirror each other in the legs as well as in the shape.
% The centre is a diamond and takes {c}, from its own border.
%
% The page reads downward in the order the parts are applied: the state, then
% the gate on the physical legs of the two sites it acts on -- the physical
% indices, not the bonds -- then the legs that come out of it, which end level
% with the three that passed the gate by.
\documentclass[border=4pt]{standalone}
\usepackage{amsmath,amssymb}
\usepackage{tikz-tensors}
\begin{document}
\begin{tikzpicture}[x=1.35cm]
  \node[canl]   (A1) at (0,0) {$A$};
  \node[canl]   (A2) at (1,0) {$A$};
  \node[centre] (C)  at (2,0) {$C$};
  \node[canr]   (B1) at (3,0) {$B$};
  \node[canr]   (B2) at (4,0) {$B$};

  % bonds, carrying the gauge direction: everything points at the centre
  \draw[bond] (A1) -- (A2);  \draw[bond] (A2) -- (C);
  \draw[bond] (B2) -- (B1);  \draw[bond] (B1) -- (C);
  \node[leg, above=1pt] at (0.5,0.25) {$D$};

  % the gate, sized and placed from the two sites it acts on
  \tngate{U}{$e^{-i\,\delta t\,\hat h}$}{C,B1}{-1.15}
  \tngatelegs{U}{north}{c}{C}
  \tngatelegs{U}{north}{r}{B1}

  % physical legs, all ending level: from the corner of each triangle, from the
  % border of the diamond, and through the gate where there is one
  \tnlegs{l}{-2.3}{A1,A2}
  \tnlegs{r}{-2.3}{B2}
  \draw[disc] (U.south -| C) -- (U.south -| C |- 0,-2.3);
  \draw[disc] (U.south -| B1.corner 2) -- (U.south -| B1.corner 2 |- 0,-2.3);

  \foreach \p/\l in {{A1.corner 3}/1, {A2.corner 3}/2, {C}/3,
                     {B1.corner 2}/4, {B2.corner 2}/5}
    \node[leg, below] at (\p |- 0,-2.3) {$n_{\l}$};
\end{tikzpicture}
\end{document}
